\chapter{Related Work}
\label{ch:related-work}

%% ── 3.1 ─────────────────────────────────────────────────────────────────────
\section{Federated Learning for IoT}
\label{sec:fl-iot}

% Survey of FL on resource-constrained IoT devices
% Ramadan et al. (2025): FL + TinyML challenges and opportunities
% Key challenges: communication overhead, device heterogeneity, data heterogeneity
% Adaptive FL: clustering, personalization, asynchronous aggregation

%% ── 3.2 ─────────────────────────────────────────────────────────────────────
\section{Federated Learning over LoRaWAN: Torres Sanchez et al.}
\label{sec:torres-sanchez}

% DETAILED REVIEW — this is our primary reference for FL methodology

Torres Sanchez et al.~\cite{torressanchez2024federated} presented the first comprehensive
study of \gls{fl} within \gls{lorawan}-constrained communication. Their framework
targets anomaly detection for industrial machinery using an optimized autoencoder. The key
aspects of their work are summarized in Table~\ref{tab:torres-comparison}.

\subsection{Experimental Setup}
% - 4 IIoT prototypes monitoring construction machinery (Manitou, Atlas Copco, etc.)
% - Private LoRaWAN → TTN → MQTT → data management agent
% - 77,798 total messages; features: battery voltage, fuel, RPM, oil temp/pressure
% - FL framework: Flower (Python), 4 virtual clients, 1 central server
% - Model: compact autoencoder, 1 hidden layer (32 neurons), Tanh, Adam, MSE loss

\subsection{Key Results}
% - FL Accuracy: 92.30%, F1: 94.77%, TNR: 90.65%, TPR: 92.93%
% - Comparable to centralized OC-SVM (92.06%) and centralized AE (93.13%)
% - Best round/epoch config: 20 epochs × 4 rounds → F1 = 95.23%
% - Optimized AE: 1.39 KB per FL round; 128 neurons: 5.52 KB; 3 layers: 44 KB

\subsection{Communication Analysis}
% Message count formula: Nm = ceil(Msize / MaxPayload) × Rd
% SF7/8: 222 bytes, SF9/10: 115 bytes, SF11/12: 51 bytes
% For 1.39 KB model: 7 messages (best) to 2233 messages (worst)
% Training airtime: 52.8 minutes optimal
% KEY INSIGHT: model size must be minimized for LoRaWAN viability

\subsection{Simulation Methodology}
% Their exact quote (use as reference for our own methodology):
% "The FL framework in this experiment simulates distributed devices functioning
%  as clients orchestrated by the Flower framework..."
% Environment: i7-1165G7, 16GB RAM, WSL + Conda
% Virtual clients: "Virtual client initialization and orchestration..."
% This confirms simulation with virtual clients is standard and accepted

%% ── 3.3 ─────────────────────────────────────────────────────────────────────
\section{ML-Based Path Loss Prediction}
\label{sec:ml-pathloss-related}

% Survey of ML for path loss modeling
% Indoor vs outdoor, different frequency bands
% Common models: ANN, Random Forest, CNN for spatial data
% Environment-driven approaches: using weather/occupancy data
% Gap: no federated approach exists for path loss prediction

%% ── 3.4 ─────────────────────────────────────────────────────────────────────
\section{Research Gap and Positioning}
\label{sec:research-gap}

% CLEAR STATEMENT: No prior work combines FL + TinyML + path loss regression
% on constrained LoRaWAN devices.

\begin{table}[H]
  \centering
  \caption{Positioning of this thesis relative to Torres Sanchez et al.\ and the prior project.}
  \label{tab:torres-comparison}
  \begin{tabular}{p{3.8cm} p{3.8cm} p{3.8cm}}
    \toprule
    \textbf{Aspect} &
    \textbf{Torres Sanchez et al.} &
    \textbf{This Thesis} \\
    \midrule
    Application domain &
      Anomaly detection (machinery) &
      Path loss regression (indoor LoRaWAN) \\[4pt]
    ML task &
      Classification / reconstruction &
      Regression ($R^2$, RMSE) \\[4pt]
    Model &
      Autoencoder (32 neurons) &
      Dense NN (5$\to$8$\to$1, 57 params) \\[4pt]
    FL framework &
      Flower (Python) &
      Custom FedAvg (TensorFlow/Keras) \\[4pt]
    Clients &
      4 virtual clients &
      6 virtual clients (one per real device) \\[4pt]
    Data source &
      77,798 machinery messages &
      1,715,869 real LoRaWAN measurements \\[4pt]
    Data distribution &
      Non-IID (per machine type) &
      Non-IID (per physical room location) \\[4pt]
    Hardware prototype &
      LoRaWAN prototype mentioned &
      TFLite Micro on real MKR WAN 1310 \\[4pt]
    Prior centralized baseline &
      OC-SVM, IF, AE (centralized) &
      XGBoost $R^2 \approx 0.93$ (prior project) \\[4pt]
    Key metric &
      F1 = 94.77\%, Acc = 92.30\% &
      $R^2$ = [XX], RMSE = [XX]\,dB \\
    \bottomrule
  \end{tabular}
\end{table}

% CONCLUDE: "Torres Sanchez et al. proved FL is feasible over LoRaWAN.
%  Our prior project proved environment drives path loss.
%  This thesis bridges both: federated path loss regression on real edge hardware
%  using real measurement data."
